
Es la diferencia entre el tiempo final y el tiempo de inicio, es decir, el tiempo transcurrido entre el instante justo antes del actual periodo y justo después del mismo.

Como se puede observar en la prueba sin estrés, la duración es siempre de menos de $30$ $\mu s$ (Figura \ref{FIG:DURNOESS}).  No obstante, durante la ejecución del programa \textit{stress}, la duración llega a los $100$ $\mu s$ y, en algunos casos, sobrepasa dicho límite (Figura \ref{FIG:DURESS}). Aunque puedan parecer muchos en la gráfica junto con todas las prioridades, es algo puramente visual, dado que el número de fallos (periodos que no cumplirían la restricción temporal) que sobre pasan $120$ $\mu s$ en prioridades mayores a $80$, es menor de $20$ (Figura \ref{FIG:DUREX120}).

La mejor prioridad según su media y su desviación típica es la $90$ (Figura \ref{FIG:DISTDUR90}).

\begin{figure}[Duration sin estrés]{FIG:DURNOESS}{Valores de de tiempo medio de ejecución sin \textit{stress} de distintas prioridades, siendo todos estos menores de $30$ $\mu s$.}
    \image{0.7\textwidth}{}{duration-no-stress}
\end{figure}

\begin{figure}[Duration con estrés]{FIG:DURESS}{Valores de de tiempo medio de ejecución con \textit{stress} de distintas prioridades, a pesar de que algunas exceden la franja de tiempo deseable, no son muchas (y están bajo una situación de estrés computacional que no se va a dar en una ejecución normal).}
    \image{0.7\textwidth}{}{duration-stress}
\end{figure}

\begin{figure}[Número de duration mayor que 120]{FIG:DUREX120}{Valores de \textit{duration} de más de $120$ $\mu s$.}
    \image{0.7\textwidth}{}{duration-exceeds-120}
\end{figure}

\begin{figure}[Distribución de duration con prioridad 90]{FIG:DISTDUR90}{Distribución de \textit{duration} con prioridad $90$.}
    \image{0.7\textwidth}{}{duration-distribution-priority-90}
\end{figure}

